% Check-your-understanding questions for sections/03_mono_playfair.tex and 04_hill_polyalphabetic.tex.
% Numerical answers recomputed by script (scratchpad/r2/verify2.py). Long practice set: questions/99_practice.tex (B--F).
\section{Monoalphabetic, Playfair, Hill and Polyalphabetic Ciphers}\label{q:03_subst}
\noindent{\small Notes: \S\ref{sec:monoalpha}, \S\ref{sec:playfair}, \S\ref{sec:hill}, \S\ref{sec:poly}. More problems: Practice Set B--F (\S\ref{sec:practice}).}

\subsection*{Check your understanding}
\begin{enumerate}[leftmargin=1.6em,itemsep=3pt]
  \item \textbf{The monoalphabetic cipher has $26! \approx 4 \times 10^{26}$ keys. Why is it still easy to break?} The attacker never searches the keys. Each letter always has the same substitute, so the ciphertext letter $\pi(x)$ occurs exactly as often as the plaintext letter $x$: the frequency histogram is preserved with its labels permuted. Matching the most frequent ciphertext letters to \ttx{e}, \ttx{t}, \ttx{a}, \dots\ and checking common digrams and trigrams (\ttx{th}, \ttx{the}) recovers the key letter by letter.
  \item \textbf{What are homophones, and why do they not make a monoalphabetic cipher secure?} Homophones are several ciphertext symbols for one plaintext letter, chosen at random, with the number of symbols proportional to the letter's frequency (about 13 for \ttx{e}), so that every symbol occurs about equally often. They flatten single-letter frequencies, but each plaintext letter still affects only one symbol, so digram and trigram patterns (\ttx{t} followed by \ttx{h}) survive and the cipher can still be broken.
  \item \textbf{Differentiate monoalphabetic and polyalphabetic substitution.} In a monoalphabetic cipher the relationship between a plaintext letter and its substitute is one-to-one and fixed for the whole message (Caesar, monoalphabetic table), so letter frequencies are preserved. In a polyalphabetic cipher a key selects one of several substitution rules at each position, so each occurrence of a letter may get a different substitute (one-to-many), which flattens letter frequencies (Vigen\`ere, autokey, one-time pad).
  \item \textbf{Why does Playfair split a doubled letter inside a pair with a filler, and why do I and J share a cell?} Two equal letters occupy the same cell, which is in the same row and the same column at once, so the encryption rules (which need two different cells) would conflict; splitting also hides doubled letters, which would otherwise show up in the ciphertext. I and J share a cell because the $5 \times 5$ square has only 25 cells, and J is rare and easily restored from context.
  \item \textbf{Why is the Playfair rectangle rule used unchanged for decryption?} Letters at $(r_1, c_1)$ and $(r_2, c_2)$ are replaced by $(r_1, c_2)$ and $(r_2, c_1)$, the other two corners of the same rectangle. Applying the rule to these gives $(r_1, c_1)$ and $(r_2, c_2)$ again, so the rule is its own inverse. The row and column rules are rotations, undone by moving left and up.
  \item \textbf{With the keyword MONARCHY, encrypt ``meet''.} Digrams \ttx{me et} (the \ttx{ee} falls across two pairs, so no filler). \ttx{me} is in column 1: m$\to$C, e$\to$L. \ttx{et} is a rectangle: e$\to$K, t$\to$L. Ciphertext \textbf{\ttx{CLKL}}.
  \item \textbf{Why does Playfair hide single-letter frequencies but not digram frequencies?} A letter's substitute depends on its partner, so one letter becomes several different ciphertext letters (\ttx{e}$\to$\ttx{D, O, N} in the solved question). But a given pair always gives the same ciphertext pair, so Playfair is a monoalphabetic cipher on the 676 digrams, and digram frequencies (\ttx{th}, \ttx{he}, \ttx{in}, \dots) can be matched, given a few hundred letters of ciphertext.
  \item \textbf{When is a Hill key matrix valid? Show what goes wrong otherwise.} $K$ is invertible modulo 26 if and only if $\gcd(\det K, 26) = 1$, i.e.\ $\det K \bmod 26$ is odd and not 13, because $K^{-1} = (\det K)^{-1}\operatorname{adj}K$ needs $(\det K)^{-1}$. With $K = \begin{pmatrix}2&0\\0&1\end{pmatrix}$ ($\det = 2$), \ttx{aa} and \ttx{na} both encrypt to \ttx{AA} ($13 \cdot 2 = 26 \equiv 0$), so Bob cannot decrypt.
  \item \textbf{Find the inverse of $K = \begin{pmatrix}3&3\\2&5\end{pmatrix}$ modulo 26.} $\det K = 15 - 6 = 9$, $9^{-1} \equiv 3$ ($27 = 26 + 1$). $K^{-1} = 3\begin{pmatrix}5&-3\\-2&3\end{pmatrix} = \begin{pmatrix}15&-9\\-6&9\end{pmatrix} \equiv \mathbf{\begin{pmatrix}15&17\\20&9\end{pmatrix}}$. Check: $KK^{-1} = \begin{pmatrix}105&78\\130&79\end{pmatrix} \equiv I$.
  \item \textbf{Why does the Hill cipher resist frequency analysis but fall to a known-plaintext attack?} Every ciphertext letter of a block is a combination of all $m$ plaintext letters, so single-letter (and for $m \ge 3$ digram) statistics are smeared out. But encryption is linear: with $m$ known plaintext/ciphertext blocks stacked as rows of $P$ and $C$, $C = PK$, so $K = P^{-1}C \bmod 26$ whenever $P$ is invertible.
  \item \textbf{How does the Kasiski test find the length of a Vigen\`ere key?} A plaintext fragment that recurs at a distance that is a multiple of the key length $m$ is encrypted with the same key letters and gives a repeated ciphertext fragment. So the distances between repeated ciphertext segments (of three or more letters) are multiples of $m$, and $m$ divides their gcd. Each of the $m$ interleaved pieces is then a Caesar cipher, broken by frequency analysis.
  \item \textbf{Encrypt ``crypto'' with the Vigen\`ere key ``abc''.} Key stream $0, 1, 2, 0, 1, 2$: $c{+}0 = $ C, $r{+}1 = $ S, $y{+}2 = 26 \equiv 0 = $ A, $p{+}0 = $ P, $t{+}1 = $ U, $o{+}2 = $ Q. Ciphertext \textbf{\ttx{CSAPUQ}}.
  \item \textbf{How does the autokey system improve on Vigen\`ere, and why is Forouzan's version still breakable in 26 trials?} The key is the keyword followed by the plaintext, so it never repeats and the Kasiski test fails. But in Forouzan's version only the first subkey $k_1$ is secret: for each of its 26 values, $P_1 = C_1 - k_1$ becomes the next key, and the whole message unrolls; only the correct $k_1$ gives readable text. (In general the key is English, so it shares the plaintext's statistics.)
  \item \textbf{Why did Vernam's cipher fail in practice, and what is the fix?} The key stream came from a loop of paper tape, so it repeated. Wherever two plaintext stretches use the same key bits, $C_1 \oplus C_2 = P_1 \oplus P_2$: the key cancels, and a guessed word in one plaintext reveals the other. The fix is a truly random key as long as the message, used only once: the one-time pad.
  \item \textbf{Classify as stream or block ciphers: additive, monoalphabetic, Vigen\`ere, Playfair, Hill.} Additive, monoalphabetic and Vigen\`ere encrypt one symbol at a time, so they are stream ciphers. Playfair (pairs) and Hill (blocks of $m$) encrypt a group together, with each ciphertext symbol depending on the whole group, so they are block ciphers.
\end{enumerate}
